\documentclass{article}
\usepackage{amsmath,amssymb,mathtools}
\usepackage[margin=1in]{geometry}
\allowdisplaybreaks
\usepackage[UTF8]{ctex}
\begin{document}

\section*{在线动态定价的递归变分重构}
\begin{equation}
\begin{aligned}
&\sup_{\Pi\in\mathcal G}\int_0^t L_q^\star(\lambda;\Pi_s,\pi_s,s)\,ds\\
&\text{subject to}\quad \Pi_t=0.
\end{aligned}\tag{A115}
\end{equation}
对每个 $q\in\mathbb N_+$ 和 $t\geq0$，定理 A1 指出该问题的最优值等于 $J^\star(\lambda;q,t)$。到达率 $\lambda>0$ 固定。决策变量为函数 $\Pi$，$\pi$ 为其密度，并非独立决策变量。被积函数中的低库存最优值由同一个在线定价问题递归给出，无需在库存水平 $q$ 下另行优化。

\subsection*{容许函数与递归被积函数}
函数类 $\mathcal G$ 由函数 $\Pi$ 组成，这些函数在 $\mathbb R_+$ 上绝对连续，且密度 $\pi\in[0,1]$。其定义条件为
\[
\Pi\in\operatorname{AC}(\mathbb R_+),\qquad
\frac{d\Pi_s}{ds}=\pi_s\quad\text{for a.e. }s\in\mathbb R_+,
\qquad \pi\text{ takes values in }[0,1].
\]
左端点可变，右端点固定。原文所列拉格朗日函数的定义域为
\[
L_q^\star:(\mathbb R_+;\mathbb R,[0,1],\mathbb R_+)\longmapsto\mathbb R
\]
函数为
\begin{equation}
L_q^\star(\lambda;\Pi,\pi,s)
=\lambda\exp(\lambda\Pi)\pi\bigl(F^{-1}(1-\pi)+J^\star(q-1,s)\bigr),
\qquad q\in\mathbb N_+.
\tag{A116}
\end{equation}
由第 3 节继承的递归边界值为
\begin{equation}
\begin{aligned}
J^\star(\lambda;q,0)&=0,&&q\in\mathbb N_+,\\
J^\star(\lambda;0,t)&=0,&&t\geq0.
\end{aligned}\tag{4, boundary values}
\end{equation}
定理 A1 给出的最优解对应关系为
\[
\pi^\star(\lambda;0,\cdot)=0,\qquad
\pi^\star(\lambda;q,s)
=\frac{\partial\Pi^\star(\lambda;q,s)}{\partial s},
\quad q\in\mathbb N_+,\quad\text{for a.e. }s\in[0,t],
\]
该关系在式 (A115) 的最优解族与在线策略问题的最优策略族之间双向成立。

\subsection*{固定输入与分布假设}
第 3 节设定单一商品、长度为 $t>0$ 的有限销售期，以及初始库存 $q$，其中
\[
\mathbb N=\{0,1,2,\ldots\},\qquad
\mathbb N_+=\{1,2,\ldots\},\qquad
q\in\mathbb N,\qquad \lambda>0.
\]
变分表述也包含边界情形 $t=0$，并适用于每个 $q\in\mathbb N_+$。顾客按已知速率 $\lambda$ 的泊松过程到达。非负估值 $X_n$（$n\in\mathbb N_+$）独立同分布，具有共同的分布函数 $F$ 及有限均值。企业知道 $F$ 及独立同分布假设，但不知道各顾客的具体估值。分布支撑的下端点为 $x_L$，上端点为 $x_U$，其中 $x_U$ 可以是 $\infty$。该分布在整个支撑上具有正密度 $f$。所有随机变量与随机过程均定义在共同的概率空间 $(\Omega,\mathcal F,P)$ 上。

企业为风险中性，商品的边际成本为零，每位顾客最多购买一件。顾客的估值不低于标价时即购买。原文给出价格与购买概率之间的关系为
\[
P(\text{purchase at price }p)=1-F(p).
\]
第 3 节给出逆函数，并指出它在 $[0,1]$ 上有唯一定义：
\[
F^{-1}(\pi)=\{x\in[x_L,x_U]:F(x)=\pi\},\qquad \pi\in[0,1].
\]
所需购买概率及其对应价格满足
\[
\pi_s\in[0,1],\qquad
p_s=F^{-1}(1-\pi_s),\qquad 0\leq s\leq t.
\]
最优策略的对应关系为
\[
\pi^\star=1-F(p^\star),\qquad
p^\star=F^{-1}(1-\pi^\star).
\]
假设存在使购买概率为零的价格，并满足
\[
1-F(p_\infty)=0.
\]
定理 A1 的前提不包含 Gumbel 吸引域、von Mises 条件、单峰性或定理 A2 中的唯一性假设。

\subsection*{购买概率策略和继承的价值函数}
以下各式定义递归构造使用的策略价值。式 (A115) 不另外引入独立的优化问题。定义在
\[
\pi=\{\pi_s,0\leq s\leq t\}
\]
这一形式下的购买概率策略，以 $(\Omega,\mathcal F,P)$ 为概率空间。属于 $\mathcal W$ 的充要条件是其取值于 $[0,1]$ 并满足以下各式：
\[
\begin{aligned}
\pi_s&=\pi(\lambda;Q_s,t-s),&&0\leq s\leq t,\\
\pi(\lambda;q,\cdot)&\text{ is measurable},&&q\in\mathbb N,\\
\pi(\lambda;0,\cdot)&=0.
\end{aligned}
\]
其中 $Q_s$ 为库存过程，$N$ 为速率为 1 的标准泊松过程：
\begin{equation}
Q_s=Q_0-N\left(\lambda\int_0^s\pi(\lambda;Q_u,t-u)\,du\right),
\qquad 0\leq s\leq t.
\tag{A108}
\end{equation}
当初始库存为 $Q_0=q$ 时，相应的期望收入为
\begin{equation}
J_\pi(\lambda;q,t)
=E\left[\left.
\lambda\int_0^t\pi(\lambda;Q_s,t-s)
F^{-1}\bigl(1-\pi(\lambda;Q_s,t-s)\bigr)\,ds
\,\right|\,Q_0=q\right].
\tag{A107}
\end{equation}
继承的最优值定义为
\begin{equation}
J^\star(\lambda;q,t)=\sup_{\pi\in\mathcal W}J_\pi(\lambda;q,t).
\tag{A109}
\end{equation}

\subsection*{首次销售分解与命题 A8}
原文将首次销售时刻定义为
\[
\tau=\inf\{s\geq0:Q_s<Q_0\}.
\]
对每个 $\pi\in\mathcal W$ 和 $q\in\mathbb N$，生存概率为
\begin{equation}
P(\tau>s\mid Q_0=q)
=\exp\left(-\lambda\int_0^s\pi(q,t-u)\,du\right),
\qquad 0\leq s<t.
\tag{A110}
\end{equation}
原文的首次销售收入恒等式为
\begin{equation}
\begin{aligned}
&E\left[\left.\int_0^\tau
\pi(Q_s,t-s)F^{-1}\bigl(1-\pi(Q_s,t-s)\bigr)\,ds
\,\right|\,Q_0=q\right]\\
&\quad=\int_0^t
\exp\left(-\lambda\int_0^s\pi(q,t-u)\,du\right)
\pi(q,t-s)F^{-1}\bigl(1-\pi(q,t-s)\bigr)\,ds.
\end{aligned}\tag{A111}
\end{equation}
中间未编号的条件恒等式为
\[
E\left[\left.\int_\tau^t
\pi(Q_s,t-s)F^{-1}\bigl(1-\pi(Q_s,t-s)\bigr)\,ds
\,\right|\,\tau,Q_0=q\right]
=J_\pi(q-1,t-\tau)/\lambda.
\]
由此得到原文的后续价值恒等式
\begin{equation}
\begin{aligned}
&E\left[\left.\int_\tau^t
\pi(Q_s,t-s)F^{-1}\bigl(1-\pi(Q_s,t-s)\bigr)\,ds
\,\right|\,Q_0=q\right]\\
&\quad=\int_0^t
\exp\left(-\lambda\int_0^s\pi(q,t-u)\,du\right)
\pi(q,t-s)J_\pi(q-1,t-s)\,ds.
\end{aligned}\tag{A112}
\end{equation}
对每个 $\pi\in\mathcal W$，命题 A8 指出，函数族
$\{J_\pi(\lambda;q,\cdot),q\in\mathbb N\}$ 是下式的唯一解：
\begin{equation}
J_\pi(\lambda;q,t)
=\int_0^t L_q\bigl(\Pi(\lambda;q,s,t),\pi(q,s),s\bigr)\,ds,
\qquad t\geq0,\quad q\in\mathbb N_+,
\tag{A113}
\end{equation}
其中
\begin{equation}
\Pi(q,s,t)=-\int_s^t\pi(q,u)\,du,
\qquad s\geq0,
\tag{A114}
\end{equation}
以及
\[
\begin{gathered}
L_q:(\mathbb R_+;\mathbb R,[0,1],\mathbb R_+)\longmapsto\mathbb R,\\
L_q(\lambda;\Pi,\pi,s)
=\lambda\exp(\lambda\Pi)\pi
\bigl(F^{-1}(1-\pi)+J_\pi(q-1,s)\bigr),
\qquad q\in\mathbb N_+,\\
J_\pi(0,\cdot)=0.
\end{gathered}
\]
对每个 $\pi\in\mathcal W$、$q\in\mathbb N_+$ 和 $t\geq0$，原文指出 $\Pi(q,\cdot,t)$ 在 $\mathbb R_+$ 上绝对连续，密度为 $\pi(q,\cdot)$。由此构造出 $\mathcal G$ 及式 (A115)。式 (A116) 的后续价值项为最优值 $J^\star$，而命题 A8 中使用的是固定策略下的价值 $J_\pi$。

\paragraph{原文记法与适用范围}
式 (A110) 至 (A116) 沿用原文在不同公式中保留或省略 $\lambda$ 的记法。具体而言，式 (A113) 中的 $L_q$ 有三个自变量，而紧随其后的定义有四个自变量 $(\lambda;\Pi,\pi,s)$。式 (A114) 省略了 $\lambda$，未在 $\Pi$ 和 $\pi$ 中显示该参数；式 (A116) 中的 $J^\star(q-1,s)$ 也没有写出 $\lambda$。按定理 A1，递归项是式 (3) 在固定到达率下给出的后续价值，因此 $\lambda$ 不应被重设为 1。第 4 节另用库存与市场规模两个自变量表示价值函数，该记法不能用于补填附录 D 公式中省略的自变量。

原文引入 $\tau$ 时写作“$\pi\in\mathcal V$”，其中 $\mathcal V$ 是定价策略类，但后续式 (A110) 和命题 A8 明确使用 $\pi\in\mathcal W$。另外，式 (A111) 的积分上限是 $\tau$，并未在 $t$ 截断；其后的条件恒等式和式 (A112) 使用下限 $\tau$，且未附加 $\tau\leq t$ 的限制。按 $\tau$ 的定义，有限销售期结束前不保证发生销售。中间的说明也未为后续价值中的 $q-1$ 新增库存为正的限定。这里保留原文表达式及其适用范围，不添加截断、指示函数，也不将定义域扩展到负剩余时间或负库存。命题 A8 本身明确将递归限制在 $q\in\mathbb N_+$。

\subsection*{后续最优值的原始定义}
定理 A1 明确调用式 (3)。下列原始定价策略公式用于阐明其含义。式 (A115) 只含一次定价优化。策略类 $\mathcal V$ 由取值于 $\mathbb R_+\cup p_\infty$ 的策略
\[
p=\{p_s,0\leq s\leq t\}
\]
组成，定义在 $(\Omega,\mathcal F,P)$ 上，满足
\[
\begin{aligned}
p_s&=p(\lambda;Q_s,t-s),&&0\leq s\leq t,\\
p(\lambda;q,\cdot)&\text{ is measurable},&&q\in\mathbb N,\\
p(\lambda;0,\cdot)&=p_\infty.
\end{aligned}
\]
其状态和收入分别定义为
\begin{equation}
Q_s=Q_0-N\left(\lambda\int_0^s
\bigl(1-F(p(\lambda;Q_u,t-u))\bigr)\,du\right),
\qquad 0\leq s\leq t,
\tag{2}
\end{equation}
\begin{equation}
J_p(\lambda;q,t)
=E\left[\left.\lambda\int_0^t
p(\lambda;Q_s,t-s)
\bigl(1-F(p(\lambda;Q_s,t-s))\bigr)\,ds
\,\right|\,Q_0=q\right],
\tag{1}
\end{equation}
以及
\begin{equation}
J^\star(\lambda;q,t)=\sup_{p\in\mathcal V}J_p(\lambda;q,t).
\tag{3}
\end{equation}
利用对应关系 $\pi=1-F(p)$ 得到式 (A107) 至 (A109)。式 (1) 至 (3) 及式 (A107) 至 (A114) 依次定义了所选变分模型中的后续价值，不引入额外的控制或约束。

\paragraph{端点表达式}
原文允许 $\pi=0$ 属于 $\mathcal G$，将 $F^{-1}$ 定义在 $[0,1]$ 上，且允许 $x_U=\infty$。因此，式 (A107)、命题 A8 及式 (A116) 均保留原文因子 $\pi F^{-1}(1-\pi)$。当分布上端点为无穷且 $\pi=0$ 时，所引原文未明确规定该乘积的处理方式，这里也不增加约定。

\end{document}
